\begin{abstract}
% Power electronic converters (PECs) are crucial in modern power systems, particularly in the integration of renewable energy sources, electrified transportation, microgrids, and portable electronics. While PECs are known for their high energy efficiency and controllability, they encounter persistent challenges, such as nonlinear dynamics, the need for real-time control, multi-objective design trade-offs, and uncertainties in operating conditions. These challenges often lead to complex, high-dimensional, and non-convex optimization problems, which are difficult to solve using traditional model-based or rule-based techniques. Deep reinforcement learning (DRL), a model-free framework combining deep learning and reinforcement learning, emerges as a promising solution for sequential decision-making in complex environments. This paper provides a comprehensive review of DRL applications in PECs, encompassing design optimization, intelligent control, and predictive maintenance. By synthesizing state-of-the-art developments, the review compares algorithmic frameworks, learning strategies, and hardware deployment practices, revealing the strengths, limitations, and unresolved issues of existing approaches. Finally, it outlines future research opportunities toward intelligent, adaptive, and self-optimizing power electronic converters capable of real-time learning and fault-tolerant operation.
% 精简版
Power electronic converters (PECs) are indispensable to modern power systems yet remain challenging to design and control because of nonlinear dynamics, real-time constraints, and multi-objective operating requirements. 
Deep reinforcement learning (DRL), a model-free framework combining deep learning with reinforcement learning, has recently emerged as a viable paradigm for sequential decision-making in complex converter environments. 
This review provides a systematic and comprehensive synthesis of DRL techniques applied to PECs across three key life-cycle stages: design optimization, intelligent control, and predictive maintenance. 
It quantifies literature trends, compares representative algorithms, and analyzes deployment practices to clarify mechanisms, benefits, and limitations. 
The discussion emphasizes physics-informed reward shaping, hybrid model–learning pipelines, and embedded implementations as key enablers of reliable and certifiable converter control. 
Finally, open challenges and research frontiers are articulated for hardware-centric validation, safety-guaranteed learning, and scalable multi-agent coordination toward fully adaptive, self-optimizing power electronic systems.

\end{abstract}
